% !TEX root = ../main.tex
\chapter{耦合模理论、微扰观点与模式分裂}
\label{ch:cmt}

\section*{学习目标}
\begin{itemize}
  \item 区分时间耦合模理论（Temporal Coupled-Mode Theory, TCMT）、空间耦合模理论与 PCSEL 专用 CWT 的角色差异。
  \item 掌握微扰公式如何把“小几何变化”“小折射率漂移”“小损耗改动”转换为频率漂移和损耗变化。
  \item 理解模式分裂不是数学细节，而是偏振选择、模竞争和远场改变的直接来源。
\end{itemize}

\section*{先修知识}
建议已掌握 \cref{ch:cwt-min,ch:cwt-3d,ch:polarization} 中关于 CWT、开放系统和对称性破缺的基础内容。

\section*{本章路线图}
\ChapterModelLevel{L1→L2}
\cref{ch:cwt-min,ch:cwt-3d} 给出 PCSEL 专用的耦合波语言，适合描述晶格散射网络。本章引入更通用的耦合模语言，先说明 TCMT 与 CWT 的关系，再用微扰公式描述结构小改动对频率和损耗的影响，最后将这些公式用于模式分裂与对称性破缺。

\section{为什么在 CWT 之外还要学耦合模理论}
很多读者看到这里会自然提问：既然前两章已经有了 CWT，为什么还要再学一个耦合模理论？答案在于两者关注的对象不同。

CWT 的自然基底是由倒格矢连接的内部平面波分量，用于解释光子晶体如何组织反馈和辐射。TCMT 则把结构压缩为一个或几个与端口耦合的共振模；其基底是已形成的共振超模及其输入输出通道。

因此，当问题关注孔形如何改变 $x$ 向波和 $y$ 向波之间的散射时，CWT 通常更自然；当问题关注某个共振模与上方自由空间、下方衬底或波导端口的耦合时，TCMT 更简洁。两者不是替代关系，而是对同一系统的不同投影。

\begin{IntuitionBox}
CWT 描述腔内场如何由周期散射形成；TCMT 描述已形成的腔模如何与外界端口交换能量。前者强调内部结构，后者强调端口与通道。PCSEL 同时具有周期开放电磁系统和可测量面发射器件两种属性，因此两者都需要。
\end{IntuitionBox}

\section{时间耦合模理论的基本语言}
先考虑最简单的单模情形。若一个共振模的复振幅为 $a(t)$，其与外部输入输出端口耦合，则 TCMT 的标准写法是
\begin{align}
\dot a &= \left(\ii\omega_0-\gamma_\mathrm{i}-\gamma_\mathrm{e}\right)a + \bm{k}^{\mathsf T}\bm{s}_+,
\label{eq:tcmt_single_a}\\
\bm{s}_- &= \mathbf{C}\bm{s}_+ + \bm{d}\,a.
\label{eq:tcmt_single_s}
\end{align}
这里，$\omega_0$ 是共振频率，$\gamma_\mathrm{i}$ 是内部损耗，$\gamma_\mathrm{e}$ 是向外部端口泄漏的外耦合损耗，$\bm{s}_+$ 和 $\bm{s}_-$ 分别表示输入与输出通道振幅，$\mathbf{C}$ 是直接散射背景矩阵，$\bm{k}$ 与 $\bm{d}$ 表示模式与端口之间的耦合系数。

为了让这组式子真正能和 RCWA、FDTD/FEM 的输出对接，还必须先说明归一化。教材里最稳妥的约定是把 $a(t)$ 视为\textbf{能量振幅}，于是
\begin{equation}
|a|^2 \propto U_{\mathrm{cav}},
\qquad
|s_\pm|^2 \propto P_{\mathrm{port}},
\label{eq:tcmt_norm_units}
\end{equation}
其中 $U_{\mathrm{cav}}$ 是腔内储能，$P_{\mathrm{port}}$ 是端口功率通量。于是 $a$ 的量纲更接近 $\sqrt{\mathrm{J}}$，$s_\pm$ 的量纲更接近 $\sqrt{\mathrm{W}}$，而 $\bm{k}$、$\bm{d}$ 的量纲更接近 $\mathrm{s}^{-1/2}$。若换用别的归一化，例如把 $|a|^2$ 直接定义为功率，那么耦合系数的量纲和数值都会跟着变。跨章节、跨软件比较 TCMT 参数时，若不先统一这一步，单看数字大小往往没有意义。

实际从数值结果提取 TCMT 参数时，最常用的接口是被动反射/透射谱的共振峰拟合。以单端口近似为例，共振附近的反射系数可写成
\begin{equation}
r(\omega)\approx C_{11}+\frac{\mathrm{i}d_1k_1}{\omega-\omega_0+\mathrm{i}(\gamma_i+\gamma_e)},
\label{eq:tcmt_r_fit}
\end{equation}
其中极点线宽由 $\gamma_i+\gamma_e$ 决定（无显著背景干涉时对应半高全宽 $2(\gamma_i+\gamma_e)$），非共振背景由 $C_{11}$ 给出。通过对 RCWA 或 FDTD 计算的谱线做拟合，可提取 $\gamma_e$（外耦合损耗）与 $\gamma_i$（内部损耗）之和；再结合被动 $Q$（$Q=\omega_0/2(\gamma_i+\gamma_e)$）和上下出光比，可进一步分离 $\gamma_e$ 的上下通道贡献。

这组式子的物理意义如下：第一式表示模幅按自身复频率演化并受外部输入驱动；第二式表示输出同时包含非共振直接散射背景和经由共振模再辐射的部分。

\begin{figure}[htbp]
  \centering
  \begin{tikzpicture}[x=1cm,y=1cm,>=Latex]
    \draw[rounded corners,fill=blue!10,thick] (2.6,0.6) rectangle (5.4,3.0);
    \node at (4.0,1.85) {腔模 $a(t)$};
    \draw[->,thick] (0.2,2.3)--(2.6,2.3);
    \node[left] at (0.2,2.3) {$s_{+,1}$};
    \draw[->,thick] (7.8,2.3)--(5.4,2.3);
    \node[right] at (7.8,2.3) {$s_{+,2}$};
    \draw[->,thick] (2.6,1.3)--(0.2,1.3);
    \node[left] at (0.2,1.3) {$s_{-,1}$};
    \draw[->,thick] (5.4,1.3)--(7.8,1.3);
    \node[right] at (7.8,1.3) {$s_{-,2}$};
    \draw[dashed,gray!70] (1.2,2.75)--(1.2,0.85);
    \draw[dashed,gray!70] (6.8,2.75)--(6.8,0.85);
    \node[font=\small] at (1.2,3.15){端口 1};
    \node[font=\small] at (6.8,3.15){端口 2};
    \node[font=\small] at (4.0,0.25){直接散射由 $\mathbf C$ 描述，共振再辐射由 $\bm d a$ 描述};
  \end{tikzpicture}
  \caption{TCMT 的最小“端口--腔模--端口”图像。PCSEL 语境下，端口可以是上方自由空间、下方衬底辐射通道，或等效的外部耦合通道；TCMT 正是把内部共振模与这些外部端口之间的能量交换压缩成有限维输入输出模型。}
  \label{fig:tcmt_port_cavity}
\end{figure}

对 PCSEL 来说，若我们把“上方自由空间出射”“下方衬底辐射”甚至“侧向端口”都视为外部通道，那么 TCMT 很适合描述谱线、共振峰、临界耦合和端口干涉现象。它尤其适用于分析被动散射谱和端口响应，而不要求我们显式追踪内部每一个倒格波分量。

\section{从单模到多模：模式竞争和分裂如何写进 TCMT}
PCSEL 更常见的是多模或近简并模情形。此时把单个振幅 $a$ 升级成模矢量 $\bm{a}$，就得到
\begin{align}
\dot{\bm{a}}
&=
\left(\ii\mathbf{\Omega}-\mathbf{\Gamma}\right)\bm{a}
+
\mathbf{K}\bm{s}_+,
\label{eq:tcmt_multi_a}\\
\bm{s}_-
&=
\mathbf{C}\bm{s}_+ + \mathbf{D}\bm{a}.
\label{eq:tcmt_multi_s}
\end{align}
其中 $\mathbf{\Omega}$ 是模间保守耦合矩阵，$\mathbf{\Gamma}$ 是损耗与辐射矩阵。到这里你会发现，这和前两章的 CWT 形式已经非常接近。区别不在数学结构，而在基底和参数来源：
\begin{itemize}
  \item 在 CWT 中，矩阵索引更偏向内部平面波或“平面波 $\times$ 纵模”基底；
  \item 在 TCMT 中，矩阵索引更偏向已经形成的共振超模以及它们与端口的关系。
\end{itemize}

因此，TCMT 与 CWT 可视为同一有限维投影思想在不同基底上的表述。实际工作中，可先用 CWT 识别重要内部耦合通道，再把候选超模参数压缩为 TCMT 参数，用于解释散射谱、反射峰、Fano 线形和端口干涉。

\section{微扰观点为什么如此重要}
在研究和设计阶段，常见问题通常是小改动：孔半径轻微变化、刻蚀深度偏差、温度升高导致折射率上升、重掺杂层吸收增加、金属位置偏移等。若每次都依赖全波重算，代价较高且不利于分析机制。微扰理论的作用，是把这些小改动映射为本征频率和损耗的一阶变化。

对于封闭或近似厄米系统，最熟悉的一阶频移公式是
\begin{equation}
\frac{\delta\omega}{\omega}
\approx
-
\frac{1}{2}
\frac{\int_V \delta\varepsilon(\rvec)\,|\E(\rvec)|^2\,\dd V}
{\int_V \varepsilon(\rvec)\,|\E(\rvec)|^2\,\dd V}.
\label{eq:perturb_hermitian}
\end{equation}
这个公式的意义并不难：如果某个区域的介电常数增加，而该模在这个区域正好场强很强，那么模式更“愿意”把能量放进去，频率通常会向低频方向漂移。反之，若扰动发生在场很弱的区域，频移就会小得多。

\cref{eq:perturb_hermitian} 的价值在于把设计动作改写为场重叠问题。这与前文 CWT 中耦合系数的结构一致：扰动是否有效，取决于它是否作用在目标模式具有显著场强或双正交重叠的区域。

\begin{RigorousBox}
\cref{eq:perturb_hermitian} 严格依赖厄米归一化和小扰动条件。对强开放的 PCSEL 辐射模，真正严格的处理需要使用准正规模（quasinormal mode）或等价的开放系统归一化框架。工程上常把 \cref{eq:perturb_hermitian} 当作第一近似使用，但必须知道它在强辐射、强非厄米情况下只是一种近似，不是普适精确公式。
\end{RigorousBox}

\begin{NoteBox}{色散材料的微扰修正}
公式\ref{eq:perturb_hermitian}假设非色散材料 ($\partial\varepsilon/\partial\omega \approx 0$)。对色散材料，严格的一阶频移公式需包含色散修正项~\cite{sauvan2013qnm,alpeggiani2017qnm}:
\begin{equation}
\frac{\delta\omega}{\omega} \approx -\frac{1}{2} \frac{\int_V \left[\delta\varepsilon + \omega\frac{\partial(\delta\varepsilon)}{\partial\omega}\right]|\E|^2\dd V}{\int_V \left[\varepsilon + \omega\frac{\partial\varepsilon}{\partial\omega}\right]|\E|^2\dd V}.
\label{eq:perturb_dispersion}
\end{equation}
在 PCSEL 常见的工作波长附近 (近红外),GaAs 和 InP 的材料色散 $\partial n/\partial\omega$ 贡献通常小于 10\%。但对高精度建模 (如波长调谐 $>10\,\mathrm{nm}$) 或宽光谱应用，建议纳入色散修正以避免系统性偏差。
\end{NoteBox}

\section{先用一个 \texorpdfstring{$2\times2$}{2x2} 非厄米玩具模型把左模/右模讲明白}
在直接写开放 Maxwell 的一般公式之前，先看一个最小非厄米矩阵
\begin{equation}
\mathbf{M}=
\begin{bmatrix}
a & b\\
c & d
\end{bmatrix},
\qquad b\neq c\ \text{或}\ a,d,b,c\in\mathbb{C}.
\label{eq:nh_2x2_matrix}
\end{equation}
设其某个本征值为 $\lambda$。对应的右本征矢满足
\[
\mathbf M \bm v^{\mathrm R}=\lambda \bm v^{\mathrm R},
\]
可取
\begin{equation}
\bm v^{\mathrm R}=
\begin{bmatrix}
b\\
\lambda-a
\end{bmatrix}.
\label{eq:nh_2x2_right}
\end{equation}
左本征矢则满足
\[
(\bm v^{\mathrm L})^{\mathsf T}\mathbf M=\lambda (\bm v^{\mathrm L})^{\mathsf T},
\]
等价于 $\mathbf M^{\mathsf T}\bm v^{\mathrm L}=\lambda \bm v^{\mathrm L}$，因此可取
\begin{equation}
\bm v^{\mathrm L}=
\begin{bmatrix}
c\\
\lambda-a
\end{bmatrix}.
\label{eq:nh_2x2_left}
\end{equation}
对比 \cref{eq:nh_2x2_right,eq:nh_2x2_left} 可见：除非问题退回厄米或特定 complex-symmetric 情形，左矢量并不等于右矢量的复共轭。在非厄米问题中，关闭投影依赖左/右本征矢的双线性配对，而不是右本征矢自身的厄米内积。

这就是开放 Maxwell 微扰中引入伴随模的最小代数原因。左模可先理解为投影时的测试函数；在开放系统中，$\E^{\mathrm L}$ 与 $\E^{\mathrm R}$ 分别由伴随 Maxwell 算子和原始 Maxwell 算子给出。
\begin{ExampleBox}{数值演示：$2\times2$非厄米矩阵的左右模}
取$a=1, d=2, b=0.1, c=0.3$，则矩阵$\mathbf{M}=\begin{bsmallmatrix}1&0.1\\0.3&2\end{bsmallmatrix}$。本征值为
\[
\lambda_\pm = \frac{3\pm\sqrt{1+0.12}}{2} \approx \frac{3\pm1.058}{2}.
\]
对$\lambda_+ \approx 2.029$，右模为$\bm{v}_+^R=[0.1,\lambda_+-1]^T \approx [0.1,1.029]^T$，
左模需解$\bm{v}_+^L{}^T\mathbf{M}=\lambda_+\bm{v}_+^L{}^T$，得到$\bm{v}_+^L=[0.3,\lambda_+-1]^T \approx [0.3,1.029]^T$。
验证双正交：$\bm{v}_-^L{}\cdot\bm{v}_+^R = 0.3\times0.1 + (\lambda_--1)(\lambda_+-1)=0$（利用韦达定理可严格证明）。
这个简单例子说明：即使在$2\times2$情况下，左模也不等于右模的复共轭或转置。
\end{ExampleBox}

\begin{ApplicationBox}{用非厄米微扰公式估计孔形误差的影响}
设设计好的圆孔 PCSEL 因工艺误差变成椭圆，介电常数扰动为$\delta\varepsilon(\bm{r})$。
若已知理想结构的右模场分布$\bm{E}_n^R(\bm{r})$和左模分布$\bm{E}_n^L(\bm{r})$，
则频率移动的一阶近似为
\[
\delta\omega_n \approx -\frac{\omega_n}{2}\frac{\int\delta\varepsilon\,\bm{E}_n^L\cdot\bm{E}_n^R\dd V}
                                      {\int\varepsilon\,\bm{E}_n^L\cdot\bm{E}_n^R\dd V}.
\]
注意分子中出现的是左模与右模的内积，而不是通常的$|\bm{E}|^2$。如果错误地使用右模自身做投影，会得到错误的频移预测，尤其在强辐射损耗模式下误差更大。
\end{ApplicationBox}
\section{为什么非厄米一阶扰动不能直接照搬厄米公式}
上一节只给出了代数例子。开放系统微扰的关键点是：\textbf{非厄米问题中，右模本身不足以关闭微扰公式。} 原因不只是频率变为复数，而是开放 Maxwell 算子不再自伴，右模之间也不再按厄米内积正交。于是，厄米情形默认成立的三件事会失效：
\begin{enumerate}
  \item 归一化不再能简单用 $\int \varepsilon |\E|^2 \dd V=1$ 关闭；
  \item 投影时不能默认用右模自己的复共轭来消去一阶模形修正；
  \item 频移与损耗修正必须被统一看成复频率修正，而不是先算实部、再临时补虚部。
\end{enumerate}

最稳妥的写法，是先把开放系统本征问题看成一般的非线性非厄米本征问题：
\begin{equation}
\hat{\mathcal L}(\tilde{\omega},p)\E^{\mathrm{R}} = 0,
\label{eq:open_right_mode}
\end{equation}
其中 $p$ 表示几何、材料或工作点参数。与之配套，需要引入\textbf{伴随算子}与\textbf{左模}：
\begin{equation}
\hat{\mathcal L}^{\mathrm{A}}(\tilde{\omega},p)\E^{\mathrm{L}} = 0,
\label{eq:open_left_mode}
\end{equation}
这里 $\hat{\mathcal L}^{\mathrm{A}}$ 是伴随算子，它由“内积里的分部积分”定义：
\begin{equation}
\langle u,\hat{\mathcal L}v\rangle = \langle \hat{\mathcal L}^{\mathrm{A}}u, v\rangle + \text{边界项}。
\label{eq:adjoint_def}
\end{equation}
直观上，\textbf{把算子从右边搬到左边，在你选定的内积下整理分部积分项，就得到伴随算子}。在封闭厄米问题里边界项为零且 $\hat{\mathcal L}^{\mathrm{A}}=\hat{\mathcal L}$，于是左模等于右模；在开放问题里边界项不再消失，左模与右模必须区分。

对一般非厄米问题，真正起投影作用的是左模与右模形成的双正交结构，而不是右模与自身的厄米内积。可把左模先理解为“做投影时的测试函数”，这会比“寻找另一支物理场”更直观。
现在对 \cref{eq:open_right_mode} 做一阶微扰。若参数改变 $p\to p+\delta p$，则
\begin{equation}
\left[
\left.\frac{\partial \hat{\mathcal L}}{\partial \omega}\right|_n \delta\tilde{\omega}
+ \delta \hat{\mathcal L}
\right]\E_n^{\mathrm{R}}
+
\hat{\mathcal L}_n\,\delta \E_n^{\mathrm{R}}
=
0.
\label{eq:open_perturb_expand}
\end{equation}
再用左模对它做投影，由于 $\E_n^{\mathrm{L}}$ 满足伴随本征方程，最后一项被消去，于是得到开放系统中的一阶公式
\begin{equation}
\delta\tilde{\omega}_n
=
-
\frac{\left\langle \E_n^{\mathrm{L}},\,\delta \hat{\mathcal L}\,\E_n^{\mathrm{R}}\right\rangle}
{\left\langle \E_n^{\mathrm{L}},\,\left.\dfrac{\partial \hat{\mathcal L}}{\partial \omega}\right|_n \E_n^{\mathrm{R}}\right\rangle}.
\label{eq:open_perturb_general}
\end{equation}
这是厄米频移公式在开放系统中的对应形式。分子仍表示扰动与模式的重叠，分母则不再是普通储能积分，而是\textbf{由伴随投影和频率导数共同定义的广义归一化量}。

\begin{RigorousBox}
\cref{eq:open_perturb_general} 就是本章后文和练习题默认调用的开放系统一阶微扰总公式。
如果读者只想带走一个“开放模微扰”的最小表达，应优先记住这一式，而不是把
\cref{eq:perturb_hermitian} 生硬推广到强开放情形。
\end{RigorousBox}

若材料存在频率色散，这个分母最好至少写出最小形式。对
\[
\hat{\mathcal L}(\omega)=
\nabla\times \mu_r^{-1}\nabla\times
-\frac{\omega^2}{c^2}\varepsilon_r(\omega)
\]
这类 Maxwell 算子，有
\begin{equation}
\left\langle \E_n^{\mathrm{L}},\,\frac{\partial \hat{\mathcal L}}{\partial \omega}\E_n^{\mathrm{R}}\right\rangle
\approx
-\frac{1}{c^2}
\int
\frac{\partial\!\big(\omega^2\varepsilon_r(\omega)\big)}{\partial \omega}
\,
\E_n^{\mathrm{L}}\cdot \E_n^{\mathrm{R}}
\,
\dd V.
\label{eq:open_perturb_dispersive_denominator}
\end{equation}
若再退回无色散近似，$\partial(\omega^2\varepsilon_r)/\partial\omega \to 2\omega\varepsilon_r$，分母才会进一步退化成更像“储能积分”的形式。也正因为如此，开放系统分母本质上不是普通厄米范数，而是包含材料色散信息的广义归一化量。

\begin{RigorousBox}
\cref{eq:open_perturb_general} 比 \cref{eq:perturb_hermitian} 更一般，因为它保留了开放系统中不能省略的两点：左/右模区分，以及算子对频率的显式依赖。只要材料色散、辐射边界或 PML 已进入模型，这两点通常都不能省略。
\end{RigorousBox}

\section{双正交、伴随模与 QNM 归一化在这里各自扮演什么角色}
到这里很容易产生另一个疑问：左模到底是不是右模的复共轭？答案是，一般并不是。对真正的厄米问题，左模与右模在适当内积下确实可以重合；但对开放系统，它们只在更特殊的条件下才会简化。对互易、无磁光非互易的很多 Maxwell 问题，算子常常在某个\textbf{无复共轭的双线性配对}下呈现 complex-symmetric 结构，此时左模更接近“转置伙伴”，而不是“共轭伙伴”。这也是为什么 QNM 归一化中常出现 $\E\cdot\E$ 而不是 $|\E|^2$。

在实际数值中，常见的一种 QNM 广义归一化可写成
\begin{equation}
\left\langle\!\left\langle \Psi_m,\Psi_n\right\rangle\!\right\rangle
=
\int_{\Omega+\mathrm{PML}}
\left[
\E_m\cdot \frac{\partial(\omega\varepsilon)}{\partial\omega}\E_n
-
\Hf_m\cdot \frac{\partial(\omega\mu)}{\partial\omega}\Hf_n
\right]\dd V,
\label{eq:qnm_norm_ch11}
\end{equation}
其中 $\Psi_n=(\E_n,\Hf_n)$，积分区域包含物理区与 PML。这个式子在正文里不要求读者立即拿去手算，但必须知道它传达的三条信息：
\begin{enumerate}
  \item 开放系统归一化通常不再使用普通能量范数；
  \item 材料色散会直接进入归一化分母，而不是只在“材料参数表”里出现；
  \item 一旦归一化换了，微扰公式的分母也必须随之更换，不能继续沿用封闭腔形式。
\end{enumerate}

更系统的函数空间直觉、QNM 定义和归一化讨论见 \cref{ch:maxwell,ch:math-appendix}。本章此处的目标，是让读者在使用微扰语言时，不会再把开放系统误当成“只是多了一个复数频率的封闭腔”。

\section{开放系统中微扰不只改频率，还会改损耗}
对 PCSEL 而言，只讨论频移是不够的，因为很多设计动作更直接改写的是损耗，而非共振频率本身。举例来说，轻微打破对称性可能几乎不改变两支模式的频率，却显著改变它们与辐射通道的耦合强度，于是损耗分裂比频率分裂更明显。

在开放系统里，更一般地可以把复频率写成
\[
\tilde\omega=\omega-\ii\gamma,
\]
因此，微扰对应复频率变化 $\delta\tilde\omega$：其实部给出频移，虚部给出衰减变化。\cref{eq:open_perturb_general} 在开放系统中输出的是统一的频移与线宽变化。模式工程不仅改变频率，也改变模式与开放通道的耦合；这解释了 \cref{ch:polarization} 中偏振选择常主要体现为损耗分裂而非频率分裂。

从设计角度看，这一结论意味着：只检查能带或共振频率分裂是不够的，还必须检查辐射损耗是否重排；后者常直接决定阈值和单模性。

\section{哪些工程公式可以继续用，哪些地方必须升级到开放系统严格形式}
到这里，读者最需要的是一条清楚的使用边界，而不是更多公式。对 PCSEL 设计，实践上可分成三层。

第一层是\textbf{可以把厄米公式当作趋势代理}的情形：模式泄露较弱、候选模远离强辐射通道重组区、扰动主要发生在腔内高场区域而不是出光边界附近、并且目标模与竞争模已经明显分离。在这种情况下，\cref{eq:perturb_hermitian} 往往能较可靠地判断“往哪个方向漂移”。  

第二层是\textbf{必须至少用左/右模公式检查一次}的情形：模式本身就是强开放模、设计变量直接作用于出光孔形或上下辐射不对称、材料色散不可忽略、或你已经开始关心虚部变化、偏振损耗分裂和 quasi-BIC 附近的辐射调控。此时，若仍只用 \cref{eq:perturb_hermitian}，最容易把“频移方向判断正确”误当成“损耗排序也正确”。  

第三层是\textbf{微扰理论本身不宜作为主模型}的情形：近简并模强混合、模式身份交换、扰动显著改变阵列边界或电流孔径、热反馈导致空间分布整体重构，或系统靠近异常点（exceptional point）等强非厄米重组区域。此时应回到完整全波或多物理场本征问题重算，而不是继续扩展一阶公式。

\begin{PitfallBox}
开放系统里最常见的误判不是“公式完全错了”，而是“用对了方向判断，却错用了定量结论”。尤其在 quasi-BIC、偏振分裂和热漂移问题上，单靠厄米频移公式往往能看出趋势，却未必能给出正确的损耗排序和阈值排序。
\end{PitfallBox}

\section{两模耦合为什么足以抓住“分裂”的核心}
尽管真实 PCSEL 可能涉及多模网络，但理解模式分裂时，一个两模模型通常已经足够说明问题。设两个近简并模式的复频率分别为 $\tilde\omega_1$ 和 $\tilde\omega_2$，耦合强度为 $\kappa_{12}$ 和 $\kappa_{21}$，则耦合矩阵为
\begin{equation}
\mathbf{M}
=
\begin{bmatrix}
\tilde\omega_1 & \kappa_{12}\\
\kappa_{21} & \tilde\omega_2
\end{bmatrix}.
\label{eq:two_mode_matrix}
\end{equation}
其本征值为
\begin{equation}
\tilde\omega_{\pm}
=
\frac{\tilde\omega_1+\tilde\omega_2}{2}
\pm
\sqrt{\kappa_{12}\kappa_{21}+\left(\frac{\tilde\omega_1-\tilde\omega_2}{2}\right)^2}.
\label{eq:two_mode_split}
\end{equation}

\begin{ReciprocityConditionBox}[互易性定理的适用条件与非互易机制]
经典 Lorentz 互易定理表述为：对各向同性、无磁光效应的线性介质，若 $(\E_1,\Hf_1)$和$(\E_2,\Hf_2)$ 是同一频率下的两组源产生的场，则
\begin{equation}
\int_V (\E_1\cdot\Jf_2 - \E_2\cdot\Jf_1)\,\dd V = \oint_S (\E_1\times\Hf_2 - \E_2\times\Hf_1)\cdot\dd\bm{S}.
\end{equation}
在无耗散的封闭系统中，这进一步使耦合矩阵满足 $\kappa_{12} = \kappa_{21}^\ast$（厄米性）；互易性本身主要给出转置对称约束。

\textbf{互易性成立的条件：}
\begin{enumerate}
  \item 介质张量对称：$\hat{\varepsilon} = \hat{\varepsilon}^{\mathsf T}$, $\hat{\mu} = \hat{\mu}^{\mathsf T}$ (无磁光/手性非互易)
  \item 线性响应：无频率转换或非线性混合
  \item 封闭系统或适当的辐射边界条件
\end{enumerate}

\textbf{开放系统中的有效非厄米耦合：}即使介质本身满足互易性，开放系统的辐射通道也会诱导有效耗散耦合，使系统矩阵呈现非厄米结构。这是因为：
\begin{itemize}
  \item 模式通过共同辐射通道发生间接耦合，耦合强度依赖于远场干涉条件
  \item 辐射损耗矩阵 $\bm{\Gamma}_{\mathrm{rad}}$ 一般不是对角的，导致有效哈密顿量非厄米
  \item 在 biorthogonal 基底下，左模与右模不构成简单的复共轭关系
\end{itemize}
这并不等同于破坏 Lorentz 互易性：在互易条件下，散射矩阵仍满足 $\mathbf S=\mathbf S^{\mathsf T}$（并可在合适基底下写作 $\kappa_{12}=\kappa_{21}$）；变化的是“保守厄米”到“开放非厄米”的动力学结构。
\end{ReciprocityConditionBox}

\cref{eq:two_mode_split} 概括了模式分裂的基本机制：当两模原本简并时，非零耦合会解除简并；若 $\tilde\omega_1$ 与 $\tilde\omega_2$ 含有不同虚部，分裂还会改写损耗排序。因此，在 PCSEL 中，分裂不仅是谱学现象，也会影响模式选择。

\section{对称性破缺怎样导致模式分裂}
现在把上述两模语言落回 \cref{ch:polarization} 的物理图像。设某个理想对称器件在 Gamma 点处支持一对简并或近简并模式，它们对应两种正交偏振或两种不同对称性表示。在完全对称时，系统对这两种模式没有偏好；一旦加入椭圆孔、双晶格扰动、上下一侧更强的辐射出口，或者热、电注入造成横向不均匀，对称性平衡就会被打破，于是出现两类可能的后果。

\begin{figure}[htbp]
  \centering
  % !TEX root = ../main.tex
% Figure: Mode splitting due to perturbation in coupled resonators
% Chapter: ch11 - CMT and Perturbation Theory for Mode Splitting
% Description: Degenerate modes split when symmetry is broken

\begin{tikzpicture}[scale=1.2]
    % Left: Unperturbed degenerate system
    \begin{scope}[shift={(-5,0)}]
      % Two identical rings (resonators)
      \draw[thick, blue] (-1,0) circle (1);
      \draw[thick, red] (1,0) circle (1);
      
      % Overlap region (coupling)
      \fill[green!20] (0,-0.3) ellipse (0.3 and 0.6);
      
      % Field distribution in ring 1
      \draw[very thick, blue, domain=-180:-90, samples=30] 
        plot({-1+cos(\x)+0.2*sin(3*\x)},{sin(\x)});
      
      % Field distribution in ring 2
      \draw[very thick, red, domain=0:90, samples=30] 
        plot({1+cos(\x)-0.2*sin(3*\x)},{sin(\x)});
      
      % Energy level diagram above
      \draw[thick] (-1.5,2) -- (-0.5,2) node[midway, above] {$\omega_0$};
      \draw[thick] (0.5,2) -- (1.5,2) node[midway, above] {$\omega_0$};
      
      \node[above, font=\bfseries] at (0,2.8) {(a) 未微扰：简并态};
      \node[below, align=center] at (0,-1.5) {两个相同谐振腔\\弱耦合，频率简并};
      
      % Coupling strength indication
      \draw[<->, thick, green!60!black] (-0.3,0) -- (0.3,0);
      \node[above, green!60!black] at (0,0.2) {\scriptsize $\kappa$};
    \end{scope}
    
    % Right: Perturbed system with splitting
    \begin{scope}[shift={(2,0)}]
      % Asymmetric rings (one slightly deformed)
      \draw[thick, blue] (-1,0) circle (1);
      \draw[thick, red, rotate around={15:(1,0)}] (1,0) circle (0.95);
      
      % Reduced overlap
      \fill[orange!20] (0.2,-0.2) ellipse (0.2 and 0.4);
      
      % Modified field patterns
      \draw[very thick, blue, domain=-180:-90, samples=30] 
        plot({-1+cos(\x)+0.15*sin(2*\x)},{sin(\x)});
      
      \draw[very thick, red, domain=-30:105, samples=30] 
        plot({1+0.95*cos(\x)-0.25*sin(2*\x)},{0.95*sin(\x)});
      
      % Energy level splitting
      \draw[thick] (-1.5,2.3) -- (-0.5,2.3) node[midway, above] {$\omega_+$};
      \draw[thick] (-1.5,1.7) -- (-0.5,1.7) node[midway, below] {$\omega_-$};
      
      \draw[thick] (0.5,2.5) -- (1.5,2.5) node[midway, above] {$\omega_0+\Delta$};
      \draw[thick] (0.5,1.5) -- (1.5,1.5) node[midway, below] {$\omega_0-\Delta$};
      
      \node[above, font=\bfseries] at (0,2.8) {(b) 微扰后：简并解除};
      \node[below, align=center] at (0,-1.5) {对称性破缺\\能级分裂$2\Delta$};
      
      % Splitting magnitude
      \draw[<->, very thick, purple] (1.7,1.5) -- (1.7,2.5);
      \node[right, purple, font=\bfseries] at (1.7,2) {$2\Delta$};
    \end{scope}
    
    % Center: Transition arrow
    \draw[->, ultra thick, orange, bend left=20] (-2,1.5) to[out=30,in=150] 
      node[above, align=center, font=\small] {施加微扰\\形变/折射率变化/应力} (0,1.5);
    
    % Bottom: Mathematical formulation
    \node[draw, align=center, font=\small] at (0,-2.5) {
      微扰哈密顿量矩阵元:\\
      $H'_{mn}=\langle\psi_m^{(0)}|\delta\varepsilon|\psi_n^{(0)}\rangle$\\[4pt]
      本征值分裂:$\Delta\omega = \frac{\omega_0}{2}\sqrt{\left(\frac{\delta\varepsilon}{\varepsilon}\right)^2 + 4\kappa^2}$
    };
    
    % Physical examples sidebar
    \node[align=left, font=\footnotesize, anchor=south east] at (-3.5,-3) {
      \textbf{实际微扰来源}:\\
      $\bullet$ 制造误差（孔径偏差$\pm 2\,\mathrm{nm}$）\\
      $\bullet$ 热梯度引起的折射率不均匀\\
      $\bullet$ 机械应力导致的双折射\\
      $\bullet$ 载流子注入的空间非均匀性
    };
  \end{tikzpicture}

  \caption{耦合谐振腔系统中微扰诱导的模式分裂现象。(a) 未微扰情况：两个完全相同的环形谐振腔通过倏逝波耦合，由于系统具有镜像对称性，两个局域模式形成简并对。(b) 施加微扰后：一侧谐振腔发生轻微形变或折射率变化，对称性破缺导致简并解除，能级分裂为$\omega_\pm=\omega_0\pm\Delta$。底部公式给出了考虑耦合作用后的精确分裂表达式。实际 PCSEL 中的微扰来源包括制造误差、热梯度、机械应力和载流子注入的非均匀性等。}
  \label{fig:ch11_mode_splitting}
\end{figure}

第一类后果是频率分裂：某一支模式的有效折射率或光程改变得更多，所以其共振频率偏移更大。

第二类后果是损耗分裂：某一支模式与辐射通道或损耗层的重叠增加得更多，所以其衰减更快或更慢。

在实际 PCSEL 设计中，这两类分裂常同时出现，只是主导程度不同。对于单模和偏振工程，关键在于它们如何改写净增益排序；分裂的重要性不在于谱线分开本身，而在于它决定哪一支模在阈值附近占优。

\begin{ExampleBox}
设方格晶格中一对原本简并的正交偏振 Gamma 点模，在引入轻微椭圆孔后出现分裂。若椭圆长轴沿 $x$ 方向，那么沿 $x$ 方向偏振的模式可能首先看到更强的保守耦合修正，于是频率偏移；与此同时，它也可能因为近场对称性改变而更强地耦合到上方辐射通道，于是损耗也改变。最终究竟哪一支模式先起振，不能只看频率谁更接近增益峰，而要看复频率与模增益共同决定的净阈值排序。
\end{ExampleBox}

这个例子说明，微扰分析和两模分裂模型可直接连接对称性破缺与器件表现。

\section{TCMT、CWT 和微扰三者怎样配合使用}
现在可以把这三套语言的分工说得更清楚。

\cref{fig:ch11_cwt_cmt_mapping}给出了CWT与TCMT两种语言之间的映射关系：CWT聚焦于晶格内部的倒格矢散射网络，适合回答"哪种散射通道主导模式形成"；TCMT 聚焦于共振模与外部端口的能量交换，适合回答"模式如何与自由空间、衬底或波导耦合"。

\begin{figure}[htbp]
  \centering
  % !TEX root = ../main.tex
% Figure content only
\begin{tikzpicture}[x=1cm,y=1cm]
  \tikzset{
    stage/.style={draw=black,thick,rounded corners,minimum width=3.2cm,minimum height=1.6cm,align=center},
    arrow/.style={->,>=Stealth,very thick}
  }

  \node[stage,fill=blue!10] (cwt) at (-4,0) {CWT 本征问题\\[3pt]$\left(\ii\bm{\Omega}-\bm{\Gamma}+\bm{G}\right)\bm v=s\bm v$};
  \node[stage,fill=orange!10] (map) at (0,0) {参数映射\\[3pt]$\omega_0,\gamma_e,\gamma_i,d$};
  \node[stage,fill=red!10] (cmt) at (4,0) {TCMT 响应\\[3pt]反射/透射/线宽};

  \draw[arrow,blue!70] (cwt) -- node[above,font=\footnotesize] {提取本征频率与损耗} (map);
  \draw[arrow,red!70] (map) -- node[above,font=\footnotesize] {代入低维散射模型} (cmt);

  \node[draw=gray!70,rounded corners,fill=gray!6,align=left,font=\footnotesize,text width=13.2cm] at (0,-2.6) {
    \textbf{工作流：}
    先由 CWT 得到目标模的复频率和模矢量，再把辐射损耗、内部损耗和远场归一化信息压缩成
    TCMT 参数。这样可以把“微观的耦合矩阵”转成“可直接和线宽、反射谱、耦合强度比较”的解析模型。
  };
\end{tikzpicture}

  \caption{CWT与TCMT的语言映射。(a)CWT 视角：关注晶格内部由倒格矢驱动的散射过程，基底是行波或驻波分量，耦合系数$\kappa$来源于介电常数分布的Fourier分量；(b)TCMT视角：关注已形成共振模与外部端口的耦合，基底是共振超模，耦合系数描述能量进出速率；(c)两者的衔接：CWT 的特征模可以作为TCMT的输入，CWT 的内部散射强度可以映射为TCMT 的有效耦合率和损耗率。}
  \label{fig:ch11_cwt_cmt_mapping}
\end{figure}

\cref{fig:ch11_cwt_cmt_mapping} 的用途是区分建模对象：解释孔形如何影响偏振时，用 CWT 追踪散射通道；解释金属厚度如何影响阈值时，用 TCMT 追踪损耗通道；同时涉及两者时，再通过映射关系连接。

如果你的问题是"哪一种倒格矢散射通道在控制Gamma 点模式的重组"，优先使用 CWT。

如果你的问题是"某个共振模式和上方自由空间、下方衬底或旁路端口怎样干涉并形成谱线"，优先使用TCMT。

如果你的问题是"结构改动很小，我想先判断频率和损耗会朝哪个方向变化"，优先使用微扰公式。

三者之间通常不是并列选择，而是按问题层级配合使用。常见路径是：先用 CWT 找到主导模和主导耦合；再在该模附近用微扰公式评估小改动；最后将候选参数区间送入 TCMT 或全波散射分析，检查端口响应和输出行为。这种串联方式比无约束参数扫描更便于解释机制。

\begin{ApplicationBox}{把 tutorial 中的 Hermitian/非 Hermitian 控制翻译成三种语言}
Noda 等人的高功率 PCSEL 路线中经常出现两个参数族：一个是由面内 $180^\circ$ 与 $90^\circ$ 衍射决定的 Hermitian 耦合，另一个是经由辐射波产生的非 Hermitian 耦合\cite{noda2023tutorial,yoshida2023}。这套说法可以直接映射到本章三种语言。

在 CWT 语言中，Hermitian 耦合主要来自介电常数 Fourier 分量把四个基本 Bloch 波彼此连接；它决定 $\Gamma$ 点附近模式的保守分裂、群速度曲率和有限尺寸包络。非 Hermitian 耦合则来自同一组 Bloch 波共享上方空气、下方衬底或背面反射器形成的辐射通道；它进入有效矩阵的虚部，改写辐射常数和损耗排序。

在 TCMT 语言中，前者更像共振模之间的内部耦合矩阵 $\mathbf{\Omega}$，后者更像通过共同端口诱导的耗散矩阵 $\mathbf{\Gamma}_{\mathrm{rad}}$。在微扰语言中，双晶格间距、孔径比例、背面 DBR 相位和 p 包层厚度都不是单纯“调频率”的旋钮，而是在同时改变复频率的实部与虚部。若只追踪频率分裂而不追踪辐射常数曲率，就会漏掉大面积单模设计真正依赖的损耗裕量。

因此，本书采用的读法是：Hermitian/非 Hermitian 控制不是另一套独立理论，而是 CWT、TCMT 与开放系统微扰在高功率 PCSEL 上的一次交汇。它的任务是把目标模与竞争模的\textbf{复频率排序}做成可设计量，而不是只把某个被动 $Q$ 或某个耦合系数推到极值。
\end{ApplicationBox}

\begin{PitfallBox}
一个常见错误是把TCMT 参数、CWT参数和全波仿真输出混在一起直接比较大小，却不检查它们的归一化约定是否一致。对开放系统而言，耦合系数、模幅和端口振幅的定义如果不统一，数值上的"差一个常数"会迅速演变成物理解释上的误判。
\end{PitfallBox}

\section{为什么本章内容会反复回到后面的器件章节}
从表面上看，微扰与模式分裂像是“纯光学细节”；但在真实 PCSEL 器件中，它们会持续被电注入和温升重新激活。温度上升引起的折射率漂移、本征增益峰漂移、自由载流子吸收上升，都会通过本章的语言转化为频移、损耗变化和模式分裂再排序。也正因为如此，后面的热-电-光耦合章节并不是另起炉灶，而是把本章的线性小扰动语言变成空间分布和工作点依赖的器件语言。

因此，本章给出的不是只适用于冷腔的数学技巧，而是一种后续章节会反复使用的分析流程：对小改动，先判断其主要改变频率、损耗还是重叠；对模式切换，先判断哪个简并被打破、哪个通道被重新加权。

\section*{本章小结}
TCMT、CWT 与微扰理论是互补的三种语言。CWT 用于追踪光子晶体内部耦合网络，TCMT 用于描述共振模与外部端口的交换，微扰理论用于把小结构或材料改动映射为频率和损耗的一阶变化。对开放 PCSEL 问题，严格的一阶扰动应写成左/右模投影形式，而不能简单照搬厄米频移公式；厄米公式在弱开放条件下仍可作为趋势代理，但在强辐射、quasi-BIC、近简并重组或显著色散区，需要使用开放系统归一化甚至完整重算。模式分裂的重要性在于它会改写偏振、远场和阈值排序。有效设计分析应根据问题在这三种语言之间切换。

\section*{练习题}
\begin{enumerate}
  \item \basicq 用自己的话解释 TCMT 与 CWT 的区别，并说明为什么它们在 PCSEL 中都不可缺。

  \item \basicq 从 \cref{eq:open_perturb_general} 的推导思路出发，解释为什么开放系统里必须引入左模或伴随模，而不能只用右模自己关闭投影。

  \item \midq 结合 \cref{eq:perturb_hermitian} 与 \cref{eq:open_perturb_general}，说明为什么“把材料改在场强最大的地方”通常最有效，但并不一定总是最优设计策略。

  \item \midq 从 \cref{eq:two_mode_split} 出发，讨论在什么情况下损耗分裂会比频率分裂更直接地影响模式选择。

  \item \hardq 设某 PCSEL 在低温下稳定单偏振工作、升温后出现偏振切换。请用“微扰 + 模式分裂”的语言给出一条可能的解释链。
\end{enumerate}

\section*{建议继续阅读}
\begin{itemize}
  \item 下一章开始进入数值方法部分。PWEM 负责给出候选模目录，本章的分裂与微扰观点则解释这些候选模在后续全波分析中的频率和损耗差异。

  \item 一般光学腔与激光模理论可参考 \cite{siegman1986,coldren2012} 中关于耦合模、模竞争和频率漂移的相关章节。

  \item TCMT 与光子晶体 slab 共振的经典文献可参考 \cite{fan2002guided,fan2003fano,suh2004tcmt}。

  \item 开放系统归一化与非厄米微扰可继续阅读 \cite{leung1994perturb,muljarov2010bw,sauvan2013qnm,kristensen2015qnm}，并与 \cref{ch:math-appendix} 对照。
\end{itemize}
